\documentclass[../../main.tex]{subfiles}

\begin{document}

\section{局部凸拓扑向量空间}\label{section:ch14-局部凸}

向量空间 $X$ 上有一类非常好的拓扑, 具有这种拓扑的 $X$ 称为局部凸拓扑向量空间, 它对我们的目标来说有两个优点: 这个类足够大, 对于赋范线性空间 $X$, 它既包括 $X$ 上的由范数诱导的强拓扑, 又包括 $X$ 上的由 $X^{*}$ 的分离点的任何子空间 $W$ 诱导的弱拓扑; 另一方面, 该类拓扑足够小, 使得对具有这种拓扑的线性空间, 若 $K$ 是不包含点 $x_{0}$ 的闭凸集, 存在过 $x_{0}$ 且不包含 $K$ 的闭超平面.

对线性空间 $X$ 中的两个向量 $u$ 和 $v$, 可表示为
\begin{align*}
x=\lambda u+(1-\lambda)v,\quad 0\leqslant\lambda\leqslant 1
\end{align*}
的向量 $x$ 称为 $u$ 和 $v$ 的凸组合. $X$ 的子集 $K$ 称为是凸的, 若它包含 $K$ 中向量的所有凸组合. 线性空间的每个线性子空间是凸的, 且赋范空间的开球与闭球也是凸的.

\begin{definition}[局部凸拓扑向量空间]\label{definition:ch14-局部凸}
局部凸拓扑向量空间是一个具有以下性质的 Hausdorff 拓扑的线性空间 $X$:
\begin{enumerate}[(i)]
\item 向量加法是连续的, 即映射 $(x,y)\mapsto x+y$ 是从 $X\times X$ 到 $X$ 的连续映射.
\item 数乘是连续的, 即映射 $(\lambda,x)\mapsto\lambda\cdot x$ 是从 $\mathbf{R}\times X$ 到 $X$ 的连续映射.
\item 该拓扑在原点有一个由凸集组成的基.
\end{enumerate}
\end{definition}

对赋范线性空间 $X$, $X^{*}$ 的子空间 $W$ 称为分离 $X$ 中的点, 若对每对 $u,v\in X$, 存在 $\psi\in W$ 使得 $\psi(u)\neq\psi(v)$. 回忆一下, 对 $X^{*}$ 的子空间 $W$ 与点 $x\in X$, $x$ 关于 $W$ 弱拓扑的邻域基由形如
\begin{align*}
\mathcal{N}_{\varepsilon,\psi_{1},\cdots,\psi_{n}}(x)=\{x'\in X\mid|\psi_{i}(x-x')|<\varepsilon,1\leqslant i\leqslant n\}
\end{align*}
的集合组成, 其中 $\varepsilon>0$ 而每个 $\psi_{i}$ 属于 $W$.

\begin{proposition}\label{proposition:ch14-局部凸例}
令 $X$ 为赋范线性空间. 则线性空间 $X$ 是关于由范数诱导的拓扑以及由 $X^{*}$ 的任何分离 $X$ 中的点的子空间 $W$ 诱导的 $W$ 弱拓扑的局部凸拓扑向量空间.
\end{proposition}

\begin{proof}
首先考虑 $X$ 具有由范数诱导的拓扑. 由于该拓扑由度量诱导, 它是 Hausdorff 的. 从范数的次可加性与正齐次性我们推出向量加法与数乘是连续的. 最后, 中心在原点的每个开球是凸的, 且这样的球族是由范数诱导的拓扑在原点的一个基.

现在令 $W$ 为 $X^{*}$ 的分离点的子空间. 为证明 $W$ 弱拓扑是 Hausdorff 的, 令 $u$ 和 $v$ 为 $X$ 中的不同向量. 由于 $W$ 分离点, 存在 $\psi\in W$ 使得 $|\psi(u)-\psi(v)|=r>0$. 则 $\{x\in X\mid|\psi(u)-\psi(x)|<r/2\}$ 与 $\{x\in X\mid|\psi(v)-\psi(x)|<r/2\}$ 分别是 $u$ 和 $v$ 的不交 $W$ 弱邻域. 为证明向量加法是连续的, 令 $x_{1}$ 和 $x_{2}$ 属于 $X$. 考虑 $x_{1}+x_{2}$ 的 $W$ 弱邻域 $\mathcal{N}_{\varepsilon,\psi_{1},\cdots,\psi_{n}}(x_{1}+x_{2})$. 则 $x_{1}$ 与 $x_{2}$ 的 $W$ 弱邻域 $\mathcal{N}_{\varepsilon/2,\psi_{1},\cdots,\psi_{n}}(x_{1})$ 与 $\mathcal{N}_{\varepsilon/2,\psi_{1},\cdots,\psi_{n}}(x_{2})$ 具有以下性质:
\begin{align*}
\text{若 }(u,v)\in\mathcal{N}_{\varepsilon/2,\psi_{1},\cdots,\psi_{n}}(x_{1})\times\mathcal{N}_{\varepsilon/2,\psi_{1},\cdots,\psi_{n}}(x_{2}),\quad\text{则 }u+v\in\mathcal{N}_{\varepsilon,\psi_{1},\cdots,\psi_{n}}(x_{1}+x_{2}).
\end{align*}
因此向量加法在 $(x_{1},x_{2})\in X\times X$ 是连续的. 类似的方法可证明数乘是连续的. 最后, 原点的基本邻域形如 $\mathcal{N}_{\varepsilon,\psi_{1},\cdots,\psi_{n}}(0)$, 而该集合是凸的, 由于每个 $\psi_{k}$ 是线性的.
\end{proof}

\begin{definition}[内部点]\label{definition:ch14-内部点}
令 $E$ 为线性空间 $X$ 的子集. 点 $x_{0}\in E$ 称为 $E$ 的内部点, 若对每个 $x\in X$, 存在某个 $\lambda_{0}>0$, 使得当 $|\lambda|\leqslant\lambda_{0}$ 时 $x_{0}+\lambda\cdot x$ 属于 $E$.
\end{definition}

\begin{proposition}\label{proposition:ch14-局部凸性质}
令 $X$ 为局部凸拓扑向量空间.
\begin{enumerate}[(i)]
\item $X$ 的子集 $\mathcal{N}$ 是开的当且仅当对每个 $x_{0}\in X$ 与 $\lambda\neq 0$, $x_{0}+\mathcal{N}$ 和 $\lambda\cdot\mathcal{N}$ 是开的.
\item $X$ 的凸子集的闭包是凸的.
\item $X$ 的开子集 $\mathcal{O}$ 的每个点是 $\mathcal{O}$ 的内部点.
\end{enumerate}
\end{proposition}

\begin{proof}
我们首先证明(i). 对 $x_{0}\in X$, 定义平移映射 $T_{x_{0}}:X\to X$ 为 $T_{x_{0}}(x)=x+x_{0}$. 则由于向量加法是连续的, $T_{x_{0}}$ 是连续的. 映射 $T_{-x_{0}}$ 也是连续的且它是 $T_{x_{0}}$ 的逆. 因此 $T_{x_{0}}$ 是 $X$ 映上 $X$ 的同胚. 因此 $\mathcal{N}$ 是开的当且仅当 $\mathcal{N}+x_{0}$ 是开的. 拓扑在数乘下的不变性的证明类似.

为证明(ii), 令 $K$ 为 $X$ 的凸子集. 固定 $\lambda\in[0,1]$. 定义映射 $\psi:X\times X\to X$ 为
\begin{align*}
\text{对所有 }u,v\in X,\quad\psi(u,v)=\lambda u+(1-\lambda)v.
\end{align*}
由于数乘与向量加法是连续的, $\psi:X\times X\to X$ 是连续的, 其中 $X\times X$ 具有乘积拓扑. 连续映射将集合的闭包映到该集合的象的闭包. 因此 $\psi(\overline{K\times K})\subseteq\overline{\psi(K\times K)}$. 然而 $\overline{K\times K}=\overline{K}\times\overline{K}$. 此外, 由于 $K$ 是凸的, $\psi(K\times K)\subseteq K$. 因此 $\psi(\overline{K}\times\overline{K})\subseteq\overline{K}$. 由于这对所有 $\lambda\in[0,1]$ 成立, $K$ 的闭包是凸的.

为证明(iii), 令 $x_{0}$ 属于 $\mathcal{O}$. 定义 $g:\mathbf{R}\times X\to X$ 为 $g(\lambda,x)=\lambda\cdot x+x_{0}$. 由于数乘是连续的, 映射 $g$ 是连续的. 但 $g(0,0)=x_{0}$ 而 $\mathcal{O}$ 是 $g(0,0)$ 的邻域. 因此存在 $0\in\mathbf{R}$ 的邻域 $\mathcal{N}_{1}$ 与邻域 $\mathcal{N}_{2}$ 使得 $g(\mathcal{N}_{1}\times\mathcal{N}_{2})\subseteq\mathcal{O}$. 选取 $\lambda_{0}>0$ 使得 $[-\lambda_{0},\lambda_{0}]\subseteq\mathcal{N}_{1}$. 则当 $|\lambda|\leqslant\lambda_{0}$ 时, $x_{0}+\lambda\cdot x$ 属于 $\mathcal{O}$.
\end{proof}

\begin{proposition}\label{proposition:ch14-连续泛函判据}
令 $X$ 为局部凸拓扑向量空间而 $\psi:X\to\mathbf{R}$ 是线性的. 则 $\psi$ 是连续的当且仅当存在原点的邻域使得 $|\psi|$ 在其上是有界的, 即存在原点的邻域 $\mathcal{N}_{0}$ 与 $M>0$, 使得
\begin{align}
\text{在 }\mathcal{N}_{0}\text{ 上},\quad|\psi|\leqslant M.\label{align:ch14-连续判据}
\end{align}
\end{proposition}

\begin{proof}
首先假定 $\psi$ 是连续的. 则它在 $x=0$ 是连续的. 由于 $\psi(0)=0$, 存在 $0$ 的邻域 $\mathcal{N}_{0}$, 使得对 $x\in\mathcal{N}_{0}$, $|\psi(x)|=|\psi(x)-\psi(0)|<1$. 因此 $|\psi|$ 在 $\mathcal{N}_{0}$ 上是有界的. 为证明逆命题, 令 $\mathcal{N}_{0}$ 为 $0$ 的邻域而 $M>0$ 使得\eqref{align:ch14-连续判据}成立. 对每个 $\lambda>0$, $\lambda\cdot\mathcal{N}_{0}$ 也是 $0$ 的邻域且在 $\lambda\cdot\mathcal{N}_{0}$ 上 $|\psi|\leqslant\lambda\cdot M$. 为证明 $\psi:X\to\mathbf{R}$ 的连续性, 令 $x_{0}$ 属于 $X$ 而 $\varepsilon>0$. 选取 $\lambda$ 使得 $\lambda\cdot M<\varepsilon$. 则 $x_{0}+\lambda\cdot\mathcal{N}_{0}$ 是 $x_{0}$ 的邻域且若 $x$ 属于 $x_{0}+\lambda\cdot\mathcal{N}_{0}$, 则 $x-x_{0}$ 属于 $\lambda\cdot\mathcal{N}_{0}$, 使得
\begin{align*}
|\psi(x)-\psi(x_{0})|=|\psi(x-x_{0})|\leqslant\lambda\cdot M<\varepsilon.
\end{align*}
\end{proof}

对无穷维赋范线性空间 $X$, 下一章的定理 9 告诉我们 $X$ 上的弱拓扑是不可度量化的. 因此在无穷维空间, 在用弱序列收敛方法证明弱拓扑性质时需要留意.

\begin{example}\label{example:ch14-vonNeumann}
(von Neumann)对每个自然数 $n$, 令 $e_{n}$ 表示 $\ell^{2}$ 中的第 $n$ 个分量是 $1$ 而其他分量是零的序列. 定义
\begin{align*}
E=\{e_{n}+n\cdot e_{m}\mid n\text{ 和 }m\text{ 是任何自然数},m>n\}.
\end{align*}
我们把证明 $0$ 是 $E$ 的关于弱拓扑的闭包点但 $E$ 中不存在弱收敛到 $0$ 的序列留作练习.
\end{example}

\textbf{注} 集合 $X$ 上的两个度量诱导相同的拓扑当且仅当关于一个度量收敛的序列关于另一个度量也收敛. 对局部凸拓扑向量空间, 事情是十分不同的. 存在具有不同的局部凸拓扑的线性空间 $X$, 关于这些拓扑序列的收敛是相同的. 序列空间 $X=\ell^{1}$ 是一个经典的例子. 空间 $X$ 是具有强拓扑与弱拓扑的局部凸拓扑向量空间, 且这些拓扑是不同的(见习题 6). 然而, Schur 的一个引理断言序列在 $\ell^{1}$ 弱收敛当且仅当它在 $\ell^{1}$ 强收敛.

\textbf{注} 拓扑向量空间定义为具有使得向量的加法与数乘是连续的 Hausdorff 拓扑的线性空间. 在缺乏局部凸性下, 这样的空间可能是十分病态的. 例如, 若 $0<p<1$, 令 $X$ 为所有使得 $|f|^{p}$ 在 $[0,1]$ 上关于 Lebesgue 测度可积的 Lebesgue 可测的扩充实值函数的线性空间. 对所有 $f,g\in X$, 定义
\begin{align*}
\rho(f,g)=\int_{[0,1]}|f-g|^{p}\,dm.
\end{align*}
则将几乎处处相等的函数等同后, $\rho$ 是诱导了 $X$ 上使得向量的加法和数乘是连续的 Hausdorff 拓扑的度量. 但 $X$ 上除了零泛函外没有连续线性泛函(见习题 52). 下一节我们证明局部凸的拓扑向量空间上存在许多连续线性泛函.

\begin{exercise}\label{exercise:ch14-42}
令 $X$ 为赋范线性空间而 $W$ 是 $X^{*}$ 的子空间. 证明 $X$ 上的 $W$ 弱拓扑是 Hausdorff 的当且仅当 $W$ 分离 $X$ 中的点.
\end{exercise}

\begin{exercise}\label{exercise:ch14-43}
令 $X$ 为赋范线性空间而 $\psi:X\to\mathbf{R}$ 是线性的. 证明 $\psi$ 关于弱拓扑是连续的当且仅当它关于强拓扑是连续的.
\end{exercise}

\begin{exercise}\label{exercise:ch14-44}
令 $X$ 为局部凸拓扑向量空间而 $\psi:X\to\mathbf{R}$ 是线性的. 证明 $\psi$ 是连续的当且仅当它在原点是连续的.
\end{exercise}

\begin{exercise}\label{exercise:ch14-45}
令 $X$ 为局部凸拓扑向量空间而 $\psi:X\to\mathbf{R}$ 是线性的. 证明 $\psi$ 是连续的当且仅当存在原点的邻域 $\mathcal{O}$ 使得 $f(\mathcal{O})\neq\mathbf{R}$.
\end{exercise}

\begin{exercise}\label{exercise:ch14-46}
令 $X$ 为赋范线性空间而 $W$ 是 $X^{*}$ 的分离点的子空间. 对任何拓扑空间 $Z$, 证明映射 $f:Z\to X$ 是连续的当且仅当对所有 $\psi\in W$, $\psi\circ f:Z\to\mathbf{R}$ 是连续的, 这里 $X$ 具有 $W$ 弱拓扑.
\end{exercise}

\begin{exercise}\label{exercise:ch14-47}
证明有限维局部凸拓扑向量空间上的拓扑由范数诱导.
\end{exercise}

\begin{exercise}\label{exercise:ch14-48}
令 $X$ 为局部凸拓扑向量空间. 证明所有线性连续泛函 $\psi:X\to\mathbf{R}$ 的线性空间 $X'$ 也有一个拓扑而使它成为局部凸拓扑向量空间, 在该拓扑上, 对每个 $x\in X$, 线性泛函 $\psi\mapsto\psi(x)$ 是连续的.
\end{exercise}

\begin{exercise}\label{exercise:ch14-49}
令 $X$ 与 $Y$ 为局部凸拓扑向量空间而 $T:X\to Y$ 是线性、一对一与映上的. 证明 $T$ 是拓扑同胚当且仅当它将 $X$ 上的拓扑在原点的基映到 $Y$ 上的拓扑在原点的基.
\end{exercise}

\begin{exercise}\label{exercise:ch14-50}
令 $X$ 为线性空间而函数 $\sigma:X\to[0,\infty)$ 有以下性质: 对所有 $u,v\in X$, (i)$\sigma(u+v)\leqslant\sigma(u)+\sigma(v)$; (ii)$\sigma(u)=0$ 当且仅当 $u=0$; (iii)$\sigma(u)=\sigma(-u)$. 定义 $\rho(u,v)=\sigma(u-v)$. 证明 $\rho$ 是 $X$ 上的度量.
\end{exercise}

\begin{exercise}\label{exercise:ch14-51}
(Nikodym)令 $X$ 为 $[0,1]$ 上的所有可测实值函数的线性空间. 对所有 $f\in X$, 定义
\begin{align*}
\sigma(f)=\int_{[0,1]}\frac{|f|}{1+|f|}.
\end{align*}
(i)将 a.e. 相等的函数等同后, 用习题 50 证明 $\rho(u,v)=\sigma(u-v)$ 定义了 $X$ 上的一个度量.
(ii)证明关于度量 $\rho$, $\{f_{n}\}\to f$ 当且仅当依测度 $\{f_{n}\}\to f$.
(iii)证明 $(X,\rho)$ 是完备度量空间.
(iv)证明映射 $(f,g)\mapsto f+g$ 是 $X\times X$ 到 $X$ 的连续映射.
(v)证明映射 $(\lambda,f)\mapsto\lambda\cdot g$ 是 $\mathbf{R}\times X$ 到 $X$ 的连续映射.
(vi)证明 $X$ 上不存在非零连续线性泛函 $\psi$. (提示: 令 $\psi:X\to\mathbf{R}$ 为线性与连续的. 证明存在 $n$, 使得只要 $f$ 是长度小于 $1/n$ 的区间的特征函数, 就有 $\psi(f)=0$. 因此对所有阶梯函数 $f$, $\psi(f)=0$.)
\end{exercise}

\begin{exercise}\label{exercise:ch14-52}
(Day)对 $0<p<1$, 令 $X$ 为 $[0,1]$ 上使得 $|f|^{p}$ 可积的所有(关于 Lebesgue 测度 $m$)可测实值函数的线性空间. 对所有 $f\in X$, 定义
\begin{align*}
\sigma(f)=\int_{[0,1]}|f|^{p}\,dm.
\end{align*}
(i)用习题 50 证明 $\rho(u,v)=\sigma(u-v)$ 定义了 $X$ 上的度量.
(ii)证明具有由 $\rho$ 确定的拓扑的线性空间 $X$ 是拓扑向量空间.
(iii)对 $X$ 中的非负函数 $f$ 与自然数 $n$, 证明存在 $[0,1]$ 的划分 $0=x_{0}<x_{1}<\cdots<x_{n}=1$, 使得对所有 $1\leqslant k\leqslant n$, $\int_{x_{k-1}}^{x_{k}}f^{p}=1/n\int_{0}^{1}f^{p}$.
(iv)对 $X$ 中的非负函数 $f$ 与自然数 $n$, 证明存在函数 $f_{1},\cdots,f_{n}$, 使得对 $1\leqslant k\leqslant n$, $\rho(f_{k},0)<1/n^{1-p}\int_{0}^{1}f^{p}$ 且
\begin{align*}
f=\sum_{k=1}^{n}1/n\cdot f_{k}.
\end{align*}
(v)证明 $X$ 上不存在连续的非零线性泛函.
\end{exercise}

\begin{exercise}\label{exercise:ch14-53}
令 $\mathcal{S}$ 为所有实数序列组成的空间, 对所有 $x=\{x_{n}\}\in\mathcal{S}$, 定义
\begin{align*}
\sigma(x)=\sum\frac{|x_{n}|}{2^{n}[1+|x_{n}|]}.
\end{align*}
证明 $\sigma$ 诱导 $\mathcal{S}$ 上的度量, 使得 $\mathcal{S}$ 是具有度量拓扑的拓扑向量空间. $\mathcal{S}$ 上的最一般的连续线性泛函是什么?
\end{exercise}

\end{document}